\section{Introduction}
\label{sec:introduction}

The Vehicle Routing Problem with Time Windows (VRPTW) is a foundational NP-hard combinatorial optimization problem underpinning modern freight logistics, last-mile parcel delivery, and sustainable distribution networks \cite{Dantzig1959, Solomon1987}. In addition to the combinatorial explosion inherent in classical routing, VRPTW enforces hard customer time-window constraints: each vehicle route must simultaneously respect vehicle capacity limits and temporal feasibility at every stop.

In commercial freight operations, the optimization objective is strictly lexicographic: dispatchers prioritize minimizing the total number of deployed vehicles ($\NV$) before minimizing the total travel distance ($\TD$) \cite{Toth2014, braysy2005vehicle}. Formally, letting $\mathcal{S}$ denote the space of feasible routing solutions, a solution $s_1 \in \mathcal{S}$ is strictly preferred over $s_2 \in \mathcal{S}$ (denoted $s_1 \prec s_2$) if and only if:
\begin{multline}
  NV(s_1) < NV(s_2) \quad \lor \\
  \quad \big(NV(s_1) = NV(s_2) \;\land\; TD(s_1) < TD(s_2)\big)
  \label{eq:lexicographic_intro}
\end{multline}

This hierarchical ordering reflects the economic structure of commercial transportation, where fixed vehicle capital acquisition, maintenance, insurance, and driver shift wages dominate marginal fuel consumption.

Adaptive Large Neighborhood Search (ALNS) \cite{Ropke2006, pisinger2007general} has long served as a leading metaheuristic paradigm for VRPTW, exploring diverse neighborhood structures via iterative ruin-and-recreate (destroy and repair) operations while strictly preserving feasibility. However, standard ALNS implementations rely on a stateless roulette-wheel or multi-armed bandit mechanism that updates selection probabilities from aggregated historical reward averages. This design exhibits two structural weaknesses:
\begin{enumerate}
    \item \textbf{Stateless Operator Selection}: Heuristic selection weights depend solely on historical scores across coarse iteration segments, ignoring the instantaneous search state, such as spatial route fragmentation, remaining temporal slack, or capacity utilization.
    \item \textbf{Myopic Plateau Stagnation}: Classical selection mechanisms struggle to coordinate multi-step operator sequences required to escape structural plateaus, frequently stalling when fleet-size reduction requires evaluating temporarily non-improving intermediate states.
\end{enumerate}

In parallel, Neural Combinatorial Optimization (NCO) has emerged as a data-driven alternative, using Deep Reinforcement Learning (DRL) models such as Pointer Networks \cite{vinyals2015pointer}, Attention Models \cite{kool2019attention}, and reinforcement learning frameworks for routing \cite{nazari2018reinforcement}. While these approaches capture spatial regularities without manual heuristic engineering, direct application of pure end-to-end neural solvers to industrial VRPTW encounters significant challenges:
\begin{enumerate}
    \item \textbf{Feasibility Degradation under Tight Time Windows}: Autoregressive decoding policies often face constraint conflicts when customer time windows are narrow, leading to high constraint violation rates or requiring surrogate penalty tunings \cite{falkner2020learning, lin2021deep}.
    \item \textbf{Objective Misalignment}: Most neural formulations optimize scalarized travel costs rather than strictly enforcing the primary-secondary fleet minimization hierarchy ($\NV \succ \TD$).
    \item \textbf{Out-of-Distribution Sensitivity}: Pure constructive policies trained on synthetic uniform instances often experience performance degradation when applied to larger instances or heterogeneous spatial topologies without extensive fine-tuning.
\end{enumerate}

To bridge these paradigms, this paper presents Tri-Level Hybrid DDQN-ALNS, a learning-augmented optimization framework that embeds state-conditioned reinforcement learning directly inside the iterative search loop of ALNS. Rather than replacing the metaheuristic with an opaque black-box neural solver, our approach maintains hard operational constraint feasibility and interpretable neighborhood operators from operations research, while replacing stateless roulette-wheel selection with a multi-level reinforcement learning control architecture. 

The primary contributions of this paper are summarized as follows:
\begin{itemize}
    \item \textbf{Hierarchical Control Architecture}: We formulate metaheuristic search management as a Hierarchical Markov Decision Process (HMDP), decoupling macro-level search exploration regimes (\emph{Plateau Controller}) from micro-level destroy-repair operator selection (\emph{Operator Controller}).
    \item \textbf{State-Conditioned Micro Selection with Entropy Gating}: We develop a Dueling Double Deep Q-Network (DDQN) operating over a 20-dimensional spatiotemporal state space, incorporating a normalized Softmax-Entropy confidence gate to adaptively balance learned neural predictions with Dirichlet/Thompson sampling priors.
    \item \textbf{Learned Acceptance Criterion (LAC)}: We introduce a supervised neural acceptance model trained on delayed downstream lookahead signals, replacing static Simulated Annealing cooling schedules with state-aware evaluation of long-term solution potential.
    \item \textbf{Rigorous Scale-Aware Benchmark across 74 Instances}: We evaluate the framework across 74 standard benchmark instances (56 Solomon-100, 12 Homberger-200, and 6 Homberger-400 instances) under strictly isolated independent cold-starts, providing an empirical characterization of solution quality, fleet reduction, and computational trade-offs supported by non-parametric statistical hypothesis testing.
\end{itemize}

The remainder of this paper is organized as follows: Section~\ref{sec:related_work} reviews related literature. Section~\ref{sec:formulation} presents the formal problem formulation. Section~\ref{sec:methodology} describes the proposed Tri-Level Hybrid DDQN-ALNS methodology. Section~\ref{sec:experiments} details the empirical benchmark results, statistical tests, and ablation analyses. Section~\ref{sec:discussion} discusses scale-aware dynamics, analytical properties, and limitations. Finally, Section~\ref{sec:conclusion} concludes the paper and outlines future directions.
